\documentclass[aspectratio=169,11pt]{beamer}

\usepackage[utf8]{inputenc}
\usepackage[T1]{fontenc}
\usepackage{amsmath,amssymb}

\setbeamertemplate{navigation symbols}{}

\title{Tutorial 11}
\author{}
\date{}

% Two completely blank slides
\newcommand{\blankslides}{%
  \begin{frame}[plain]\end{frame}%
  \begin{frame}[plain]\end{frame}%
}

\begin{document}

\begin{frame}[plain]
  \titlepage
\end{frame}

% ------------------------------------------------------------
\begin{frame}{Remark}
Recall, we can solve a max-flow problem:
\begin{itemize}
  \item in $O(nm^2)$ time using Edmond-Karp (irrespective of capacities or flow)
  \item in $O(m\,|f|)$ time using flow template aka Ford-Fulkerson Method.
  \item in $O((m+n)m\log(|f|))$ using Ford-Fulkerson algorithm
\end{itemize}
In some problems, $|f|\le Cn$ for some constant $C$, in that case we can solve in
$O(m\,|f|)=O(m\cdot(Cn))=O(mn)$ using Ford-Fulkerson Method (which in this case is faster than Edmond-Karp).
\end{frame}

% ============================================================
\section{Classics}

\begin{frame}{Bipartite Matching}
Given $n_1$ students and $n_2$ different food items in the mess, we can form a bipartite graph
$G=(V_1\cup V_2,E)$ where $(u,v)\in E$ implies that the student $u$ likes the food item $v$.

\medskip
Assign food items to students, maximizing the number of students who receive a food item they like
(you don't need to assign all food items, a student might not be assigned a food item which they don't like).

\medskip
Formally, we want a the maximal matching $M$ of $G$. A matching of $G$ is a subset of edges in $G$
such that no vertex is incident to more than 1 edge included in the matching.
\end{frame}
\blankslides

% ------------------------------------------------------------
\begin{frame}{Assignment Problem}
Given a set $X$ of students and a set $Y$ of ice cream tubs.

\medskip
A student $x\in X$ can only eat $c(x)$ scoops of ice cream. An ice cream tub $y\in Y$ can only has
$c(y)$ scoops of ice cream.

\medskip
Finally, each pair of student and ice-cream $(x,y)$ can only be matched $c(x,y)$ times
(student $x$ only wants $c(x,y)$ scoops of flavour $y$).

\medskip
Use up as many scoops as possible!

\medskip
\textbf{Hint:} Try modifying the above construction. How do we enforce $c(x)$ and $c(y)$ although?
\end{frame}
\blankslides

% ------------------------------------------------------------
\begin{frame}{Decomposition of Flows}
Given a flow network $G=(V,E)$, $c:E\to\mathbb{R}_{\ge 0}$, $s,t\in V$. Let $f,f'$ be flows on $G$ and
$g$ be flow on $G_f$ (residual graph).

\medskip
Prove that the flow $f'-f$ on $G_f$ has $\mathrm{cost}(f'-f)=\mathrm{cost}(f')-\mathrm{cost}(f)$.

\medskip
(ii) Prove that the flow $f+g$ on $G$ has $\mathrm{cost}(f+g)=\mathrm{cost}(f)+\mathrm{cost}(g)$.

\medskip
(iii) Prove that $(G_f)_g=G_{f+g}$ if we count parallel edges with same cost as equal.
\end{frame}

\begin{frame}{Decomposition of Flows (continued)}
(iv) Using the above, prove the \textbf{flow decomposition theorem}: Given a flow network $(G,c,s,t)$;
the flow can be decomposed into paths and cycles.

\medskip
That is there exists a finite set $\mathcal{P}$ of paths from $s$ to $t$ and a finite set of cycles
$\mathcal{C}$ such that each path and cycle have a positive weight $w$ and
\[
  f(e)=\sum_{P\in\mathcal{P},\,e\in P} w(P)\;+\sum_{C\in\mathcal{C},\,e\in C} w(C)
\]
\end{frame}
\blankslides

% ------------------------------------------------------------
\begin{frame}{Minimum Cost Maximum Flow}
Given a flow network $G=(V,E)$, $c:E\to\mathbb{R}_{\ge0}$, $s,t\in V$ and a cost function
$\alpha:E\to\mathbb{R}$ the cost of a flow $f$ in $g$ is:
\[
  \mathrm{cost}(f)=\sum_{e\in E} f(e)\,a(e)
\]
The minimum cost maximum flow problem is to find a flow $f$ in $G$ such that the value of the flow is
maximized, and among all possible solutions with maximum value, to find the one with minimum cost.

\medskip
Assume that $G$ doesn't have negative cycles wrt $\alpha$. We want to show that this algorithm works:
\end{frame}

\begin{frame}{Minimum Cost Maximum Flow: Algorithm}
\textbf{Min Cost Max Flow}
\begin{enumerate}
  \item Initialize $f$ by setting $f(e)\leftarrow 0$ for all edges
  \item While $t$ is reachable from $s$ in $G_f$ (residual graph):
  \item[3] \quad Let $P$ be the the shortest path from $s$ to $t$ (wrt $\alpha$)
  \item[4] \quad Let $\gamma$ be the minimum capacity among the edges of $P$ in $G_f$.
  \item[5] \quad Define $\Delta f$ by putting $\gamma>0$ on every edged in $P$.
  \item[6] \quad Let $f\leftarrow f+\Delta f$
  \item[7] return $f$
\end{enumerate}
\end{frame}

\begin{frame}{Minimum Cost Maximum Flow: Questions}
(i) Show that an iteration of this algorithm doesn't end up creating negative cycles in $G_f$.
You might want to use the flow decomposition theorem.

\medskip
(ii) Show that the returned $f$ will be the minimum cost flow. Note, we can be sure the flow is maximum
as we are following the max-flow template.
\end{frame}
\blankslides

% ------------------------------------------------------------
\begin{frame}{Min Cost Fixed Flow}
Same setup as Min Cost Max Flow, but you are given the amount of flow $a$ you need to achieve and you
need to minimize the cost among flows achieving it. That is find $\min_{|f|\ge a}\mathrm{cost}(f)$.

\medskip
\textbf{Hint:} If a push would take us above $a$, just push as much to reach $a$ and then terminate.
Why is this min cost?
\end{frame}
\blankslides

% ------------------------------------------------------------
\begin{frame}{Intersection of Min Cut}
Given a flow network where all minimum cuts have a nonempty intersection, prove that there exists an
edge that is in every minimum cut. The intersections are taken edge-wise.

\medskip
\textbf{Hint:} Consider $G_f$ for max flow $f$. Can you see two min cuts (which are unique upto relabelling)?

\medskip
\textbf{Hint:} Think about their intersection now. You'll either get a contradiction or (neighboring)
vertices $u$ and $v$ such that $u$ is always reachable from $s$ and $v$ is never reachable from $s$ in $G_f$.
\end{frame}
\blankslides

% ============================================================
\section{Multiple Routing}

\begin{frame}{Two Routes with no overlapping edges}
Gojo wants Yuji and Megumi to travel from Sendai to Tokyo. However, if Yuji or Megumi end up on the same
road, they will inevitably start racing and crash their bikes.

\medskip
So, Gojo, given the road network is a directed graph $G=(V,E)$ wants to find two routes from $s$ (Sendai)
to $t$ (Tokyo) such that none of the edges is used in both the routes.

\medskip
Give an algorithms to do so in $O(nm)$ time or report that it is not possible.
\end{frame}
\blankslides

\begin{frame}{Two Routes with no overlapping vertices}
Gojo realized that Yuji and Megumi will end up doing something stupid if they ever end up in the same
city (vertex).

\medskip
So, Gojo, given the road network is a directed graph $G=(V,E)$ wants to find two routes from $s$ (Sendai)
to $t$ (Tokyo) such that none of the vertices is used in both the routes.

\medskip
Give an algorithms to do so in $O(n(n+m))$ time or report that it is not possible.
\end{frame}
\blankslides

\begin{frame}{Multiple Routes with no overlapping vertices}
Gojo realized that this is true for all the $k$ sorcerers.

\medskip
So, Gojo, given the road network is a directed graph $G=(V,E)$ wants to find $k$ routes from $s$ (Sendai)
to $t$ (Tokyo) such that none of the vertices is used in more than 1 of these routes.

\medskip
Give an algorithms to do so in $O((n+m)nk)$ time or report that it is not possible.
\end{frame}
\blankslides

\begin{frame}{Shortest multiple routes with no overlapping vertices}
Gojo also doesn't want these people to take ages in getting back. Let $\alpha:E\to\mathbb{R}_{\ge0}$ be
the time to traverse an edge (length of road).

\medskip
So now he wants to find $k$ routes from $s$ to $t$ such that no vertex is used in more than one route and
the total length of routes is minimized.

\medskip
\textbf{Hint:} Think about min cost fixed flow!
\end{frame}
\blankslides

% ============================================================
\section{Tuple Selection Problems}

\begin{frame}{Tuple Selection Problems}
The bipartite maximum matching and assignment problems are examples of a larger class of problems which
Jeff Ericson calls ``Tuple Selection''.
\end{frame}

\begin{frame}{Tuple Selection Problem}
Given finite sets $X_1,X_2,\dots,X_d$; we want to select the largest possible set of $d$-tuples of the form
$(x_1,x_2,\dots,x_d)$ where $x_i\in X_i$ subject to:
\begin{itemize}
  \item For each $i$, $x\in X_i$ can only appear in $c(x)$ many tuples
  \item For each $i$, $x\in X_i$ and $y\in X_{i+1}$ can appear together in atmost $c(x,y)$ many tuples.
\end{itemize}
Where $c(x)$ and $c(x,y)$ are positive integers or $\infty$.

\medskip
We can see the maximum bipartite matching problem was a $d=2$, $c(x)=1$ and $c(x,y)=1$ if $xy$ is an edge
in the original graph and $0$ otherwise.

\medskip
Similarly, assignment problem is the general $d=2$ case.
\end{frame}

\begin{frame}{Tuple Selection using Flows}
Show that the tuple selection problem can be solved in $O(nm^2)$ time where
\[
  n=\sum_i |X_i| \qquad m=\prod_i |X_i|
\]
\end{frame}
\blankslides

\begin{frame}{Exam Scheduling}
You have to schedule $n$ exams in $r$ rooms on $d$ dates to be invigilated by $p$ invigilators upto to
the constraints:
\begin{itemize}
  \item Each exam has to be scheduled only once.
  \item No two exams can happen in the same room, on the same date
  \item No invigilator can invigilate more than one exam on the same day.
  \item Every invigilator has a list of days on which they are available. They can't invigilate on any other days.
  \item No invigilator can invigilate more than 3 exams.
\end{itemize}
Reduce this problem to Tuple Selection.
\end{frame}
\blankslides

\begin{frame}{Tournament Elimination}
$n$ teams are playing in a tournament. Team $i$ has won $W[i]$ games and has $M[i,j]$ matches remaining
against team $j$.

\medskip
Design an algorithm, that given a team $i$, tells if it is possible for team $i$ to win the tournament
(that is have or be tied for the highest number of wins after all remaining matches are done). We are
expecting a $O(n^6)$ algorithm.

\medskip
\textbf{Hint:} Why is this question in this problem sheet? And why in this section?
\end{frame}
\blankslides

\begin{frame}{Hogwarts Renovation Committee}
The Hogwarts School of Witchcraft and Wizardry has decided to form a committee to renovate the castle.
Exactly one professor must be selected from each of the school's academic departments.

\medskip
Some professors teach in multiple departments. However, each professor can represent only one department
on the committee.

\medskip
For example, if Professor Snape teaches both Potions and Defence Against the Dark Arts, and he is selected
as the representative for Potions, then a different professor must represent Defence Against the Dark Arts.
\end{frame}

\begin{frame}{Hogwarts Renovation Committee (continued)}
Additionally, the committee must contain an equal number of professors from each of the houses:
Hufflepuff, Ravenclaw, Gryfindor and Slytherin.

\medskip
Fortunately, the total number of departments is a multiple of 4, say $4k$.

\medskip
Design an efficient algorithm to select the committee members or tells that it is impossible.
\end{frame}
\blankslides

\begin{frame}{Election}
Consider an election between two candidates, $A$ and $B$, with $n$ voters.

\medskip
Each voter $i$ has a preferred candidate $C_i\in\{A,B\}$ and a preference penalty $p_i\in\mathbb{Z}$,
representing the cost of voting against their preference.

\medskip
Voter relationships are represented by an undirected weighted graph $G=(V,E)$, where each vertex represents
a voter and each edge $(i,j)$ has weight $w_{ij}\in\mathbb{Z}$, representing the strength of their friendship.
\end{frame}

\begin{frame}{Election (continued)}
Each voter incurs:
\begin{itemize}
  \item A penalty of $p_i$ for voting against their preferred candidate.
  \item A penalty of $w_{ij}$ if they vote differently from a friend $j$.
\end{itemize}
Design an efficient algorithm that, given $G$, the voter preferences $C_i$, and the penalties $p_i$ and
$w_{ij}$, determines the voting assignment that minimizes the total penalty and identifies the winning candidate.
\end{frame}
\blankslides

% ============================================================
\section{Graph Algorithms}

\begin{frame}{Max Weight Closure}
Given a directed graph $G=(V,E)$, a subset $S\subseteq V$ is called a closure if no edge enters $S$ from
outside it. Formally, $\forall v\in S,\ N_{\mathrm{in}}(v)\cap(V\setminus S)=\emptyset$.

\medskip
Monster's University allows students to freely choose courses for their degree, subject to two rules:
\begin{itemize}
  \item To take a course, you must also take all its prerequisites.
  \item The selected courses must form a closure.
\end{itemize}
\end{frame}

\begin{frame}{Max Weight Closure (continued)}
The prerequisite structure is represented by a directed acyclic graph $G=(C,E)$, where $(u,v)\in E$ means
that $u$ is a prerequisite for $v$.

\medskip
Each course $v\in C$ has a market value $m(v)\in\mathbb{R}$. A course may even have negative value, perhaps
because it makes you more likely to pursue a PhD, making companies wary of hiring you!

\medskip
Design an algorithm that given $G$ and $m$, finds the valid degree (a closure) that maximises the total
market value:
\[
  \max_{S\subseteq C}\ \sum_{v\in S} m(v)
\]
\end{frame}
\blankslides

% ------------------------------------------------------------
\begin{frame}{Disjoint Path Cover in DAGs}
A path cover for a directed-graph $G$ is a set of (directed) paths such that every vertex of $G$ lies on
atleast one of the paths.

\medskip
We call a path cover disjoint if every vertex lies on exactly one path.

\medskip
Given a directed acyclic graph, reduce finding the disjoint path cover to a bipartite matching problem
(and hence, a flow problem!) and solve in $O(nm)$ time.

\medskip
\textbf{Hint:} What could the matching represent? What about $a$ being matched to $b$ represents there is
a path where $a\to b$ ($b$ comes after $a$)?

\medskip
\textbf{Hint:} We claim we can cover in $k$ disjoint paths if there is a matching of size $|V|-k$ in the
flow graph we make\ldots
\end{frame}
\blankslides

% ------------------------------------------------------------
\begin{frame}{Sam I Am}
Sam is facing a temple which can be described by a $m\times n$ grid and he has the locations of all enemies
in the temple (each location can be thought of as the intersection of a row and a column in this grid).

\medskip
All of a sudden, he realizes that he can kill the enemies without entering the temple using the great
cannon ball which spits out a gigantic ball bigger than him killing anything it runs into and keeps on
rolling until it finally explodes.
\end{frame}

\begin{frame}{Sam I Am (continued)}
But the cannonball can only shoot horizontally or vertically and all the enemies along the path of that
cannon ball will be killed.

\medskip
Sam wants to know the minimum number of cannon balls and the positions from which he can shoot the
cannonballs to eliminate all enemies from outside that temple.
\end{frame}

\begin{frame}{Sam I Am (Hints)}
\textbf{Hint:} Consider the graph $G=(V,E)$ where $V=\{R_1,R_2,\dots,R_m,C_1,C_2,\dots,C_n\}$ and
$(R_i,C_j)\in E\iff$ there is an enemy at the intersection of row $i$ and column $j$. What would a valid
(not minimal, just kills all enemies) shooting scheme be equivalent to? What does a matching in $G$
corresponds to in the problem instance?

\medskip
\textbf{Hint:} In the residual graph, consider the row vertices that are unreachable from $s$ and column
vertices from where it is impossible to reach $t$. Let's name these misfits. Using the
maxFlow-minCut duality, what can be said about misfits?
\end{frame}
\blankslides

% ------------------------------------------------------------
\begin{frame}{Largest Common Rooted Tree}
A rooted tree is a directed acyclic graph, in which every vertex has exactly one incoming edge, except for
the root, which has no incoming edges Equivalently, a rooted tree consists of a root vertex, which has edges
pointing to the roots of zero or more smaller rooted trees.

\medskip
Describe a polynomial-time algorithm to compute, given two rooted trees $A$ and $B$, the largest common
rooted subtree of $A$ and $B$.
\end{frame}

\begin{frame}{Largest Common Rooted Tree (Hints)}
\textbf{Hint:} An approach would be to compute $\mathrm{dp}[u,v]$ which would be the largest common rooted
subtree between $u\in A$ (and it's descendants) and $v\in B$ (and it's descendants) where for all vertices
$u$ and $v$ using dynamic programming\ldots

\medskip
\textbf{Hint:} Sadly, this is not tht simple. If the number of children was bounded or there was some
ordering over them, it might have worked. Given arbitrary number of children, we need a mapping from
children of $u$ to those of $B$ (or shall I say, matching\ldots)
\end{frame}
\blankslides

% ============================================================
\section{Fun Problem!}

\begin{frame}{Rounding a Matrix}
Suppose we are given an array $A[1\ldots m][1\ldots n]$ of non-negative real numbers. We want to round $A$
to an integer matrix, by replacing each entry $x$ in $A$ with either $\lfloor x\rfloor$ or $\lceil x\rceil$,
without changing the sum of entries in any row or column of $A$. For example:
\[
\begin{pmatrix}
1.2 & 3.4 & 2.4\\
3.9 & 4.0 & 2.1\\
7.9 & 1.6 & 0.5
\end{pmatrix}
\mapsto
\begin{pmatrix}
1 & 4 & 2\\
4 & 4 & 2\\
8 & 1 & 1
\end{pmatrix}
\]
\end{frame}

\begin{frame}{Rounding a Matrix (Parts)}
(a) Describe and analyze an efficient algorithm that either rounds $A$ in this fashion, or reports
correctly that no such rounding is possible.

\medskip
(b) Prove that a legal rounding is possible if and only if the sum of entries in each row is an integer,
and the sum of entries in each column is an integer. In other words, prove that either your algorithm from
part (a) returns a legal rounding, or a legal rounding is obviously impossible.

\medskip
\textbf{Extra:} Suppose we are guaranteed that none of the entries in the input matrix $A$ is an integer
Describe and analyze an even faster algorithm that either rounds $A$ or reports correctly that no such
rounding is possible. Your algorithm must run in $O(mn)$ time.
\end{frame}
\blankslides

% ------------------------------------------------------------
\begin{frame}{Abusing Trader Problem}
You (unilaterally) control the price of an asset $M$. A price movement $\mathcal{M}=(m_1,m_2,\dots)$ is a
sequence of $\pm1$ such that the price/value of $M$ on day $j$ is $\sum_{i=1}^{j} m_i$.

\medskip
A trade is a tuple $(w,l,f)$ where the trader makes $f(w)$ profit if $M$ reaches $w$ before reaching $l$
and makes a loss of $f(l)$ if it reaches $l$ before reaching $w$. It has value $0$ otherwise.

\medskip
Given a set of trades $\mathcal{T}$, choose a movement $\mathcal{M}$ to maximize trader loss.\footnote{There is a
$O(n^2)$ algorithm where $n$ is the number of trades. That is much harder to find but very beautiful.
Check out the paper Lafond's ``How Brokers Can Optimally Abuse Traders'', FUN 2022.}
\end{frame}
\blankslides

\end{document}